\documentclass[11pt]{article}
\usepackage[a4paper,margin=25mm]{geometry}
\usepackage{amsmath,amssymb,amsthm,booktabs}
\newtheorem{theorem}{Theorem}
\title{The square $(3,5,3,5)$ through order 3840}
\author{}\date{}
\begin{document}\maketitle
We complete the ninth bounded search associated with Table 3 of
Fernandes--Leemans--Weiss, \emph{An Exploration of Locally Spherical
Regular Hypertopes}, fixed 16-page edition. The normal quotient method,
finite target tables and exact marked coset test are acquired from the
reviewed components. Their computational premises remain explicit.

Let $W$ have four involutory generators, with consecutive square bond
orders $3,5,3,5$ and both diagonal bond orders 2. A regular hypertope
with this diagram supplies a normal subgroup of $W$ whose index is its
type-preserving group order. This follows from its distinguished chamber
involutions and does not assume that its finite-diagram residues are
universal spherical geometries.

\begin{theorem}
Subject to the acquired normal quotient enumeration premises, there is
exactly one positive ordered kernel through index 3840. It has index
1920 and four maximal parabolics of order 60. Thus 1920 is the minimum
group order when projective finite-diagram residues are allowed. There
is no example of order 3840, and no example through that bound with all
four universal spherical residues of type $H_3$.
\end{theorem}
\begin{proof}
The complete normal-kernel inventory is
\[
\begin{array}{c|rrrrrrrrr}
\text{index}&1&2&60&120&660&960&1320&1920&3840\\
\text{multiplicity}&1&1&1&1&1&1&1&7&11.
\end{array}
\]
These are 25 kernels, including the whole group. Exact rank-two
subgroups have orders 4, 6 and 10, so a required quotient has order
divisible by 60. Only the first two kernels may be omitted by this
necessary divisibility test. For each remaining kernel we construct its
faithful marked quotient and apply the acquired tests for all 256
subgroup intersections and all 32 base-flag extensions.

Twenty-one of the 23 marked quotients fail the intersection property.
This includes all eleven quotients of order 3840. The order-960 quotient
passes the intersections but fails the actual incidence test: a specified
two-type base flag has ten completions of another type, while its
stabilizer orbit has five. Consequently it is not flag-transitive and
does not supply the required regular hypertope. The sole passing kernel
has order 1920. Its independent literal coset calculation verifies all
the flag extensions, so the acquired coset theorem proves that it gives
a thin, residually connected, flag-transitive geometry with full
type-preserving group of that order.

The four local orders of this positive quotient are all 60. Each local
rank-three marked group is a quotient of the order-120 Coxeter group
$H_3$, but not the full group. In particular this positive kernel does
not satisfy the universal-residue convention. Since every other kernel
through 3840 was rejected already under the broader convention, none
satisfies the universal-residue condition.
\end{proof}

\newpage
\paragraph{The positive witness and its interpretation.}
The four involutions act on 12 points. Their complete arrays are retained
in the literal certificate; in cycle notation they are
\[
\begin{aligned}
t_1&=(1\ 3)(2\ 4)(5\ 6)(7\ 11)(8\ 12)(9\ 10),\\
t_2&=(1\ 5)(2\ 6)(3\ 4)(7\ 8)(9\ 11)(10\ 12),\\
t_3&=(1\ 2)(3\ 4)(5\ 9)(6\ 10)(7\ 12)(8\ 11),\\
t_4&=(1\ 6)(2\ 5)(3\ 7)(4\ 8)(9\ 10)(11\ 12).
\end{aligned}
\]
The literal checker closes the generated group, checks every pair order,
forms all sixteen generated subgroups and all their intersections, and
constructs every parabolic coset as a set of permutations. Each of the
four types has 32 objects. The maximal local quotients have order 60;
the earlier spherical classification identifies the corresponding
marked geometries as projective $H_3$ residues. No named global group
structure is inferred merely from the order 1920.

\paragraph{Completeness and the retained computation.}
The earlier search already returned the 25-kernel inventory and its
single positive kernel. The current certificate keeps that execution
and its exact dependencies. To make every rejection inspectable, the
new output routine also records the compact marked actions for the
negative quotients. Its only changes to the acquired caller remove the
positive-only output filter and rename the output count. The search,
classification inputs, index bound and mathematical tests are unchanged.
The checker verifies that exact program difference.

A new bounded execution supplies all 23 actions. An actual isomorphism
taking the four generators in order identifies its positive kernel with
the earlier one. On subsequent replay the same comparison also checks
the whole recorded marked inventory, preserving multiplicities.
Every GAP-to-JSON array is read from the current computation. Independent
integer permutation closure checks all 23 verdicts, including the
order-960 incidence failure. We therefore use neither an abstract
group-name match nor an order match as the certificate of a marked
quotient.

\paragraph{Source correction and limits.}
The source's first $(3,5,3,5)$ entry has order 3840. The eleven rejected
kernels exclude that order under either local convention. With the
source's broad finite-diagram meaning, its minimum claim must instead
use 1920. If universal spherical residues were intended, the conclusion
here is only a lower bound greater than 3840; no existence or least order
beyond the bound is asserted.

The second printed order, 9720, is handled in a separate divisor-search
component. This note does not repeat or acquire the unsuccessful general
search through 9720. The original targets and claims belong to the source
authors; the LINS algorithm and classification tables are prior work.
The present contribution is the last bounded application, its explicit
negative certificates and the correction with its local-scope distinction.
It is a finite published-claim correction, not a newly solved open
conjecture. Other source obligations, whole-manuscript registration, Lean
formalization and compression measurements remain separate.
\end{document}
